% ----------------------
\subsection{Section 74 (1830)}
% ----------------------
	\label{DC74}
	
	% -----
	\subsubsection{Dating this revelation}
	% -----
	
		Woodford (HDDC 908--911) associates it with the translation of the Bible, following JS's history. Marquardt does the same and offers this comment in a footnote:
		
			\begin{quote}
				This explanation of scripture has no date even though it was recorded twice in the KRB. The 1835 D\&C also has no date as to when the explanation was given. The manuscript history when including the explanation in the narrative for January 1832 dates it to ``this period'' (Jessee, \textit{Papers of Joseph Smith}, 1:371). The approximate date of 17 February-early March 1832 is based on the revision of the NT that Smith and Rigdon were working on. They were revising John 5 on 16 February and by 20 March had proceeded to about Revelation 11.
			\end{quote}
			
		The key information for dating it to 1830 is the fact that Whitmer included it in RB1; see the historical background from the JSP below. JS could have produced it at any point during a period of about a year, from most of 1830 into the early part of 1831.
	
	% -----
	\subsubsection{Historical Background}
	% -----
		\label{DC74-HB}
		
		% --
		\paragraph{JSP}
		% --
		
			``This `explanation' clarifies a New Testament verse, \hyperref[1CORINTHIANS-7-14]{1 Corinthians 7:14}, which historically had been an important passage for justifying infant baptism.1 Like the Book of Mormon, this document rejects the need for infant baptism by explaining that little children are made clean through the atonement of Jesus Christ, without baptism. Although it is possible that questions regarding 1 Corinthians 7:14 or infant baptism prompted this explanation, the precise circumstances are unknown.
			
			``The date the document was produced is also uncertain. When John Whitmer copied the document into Revelation Book 1 in 1831, he dated it simply `1830.' However, he placed it between documents dated 6 January 1831 and 4 February 1831. The editors of the Book of Commandments did not include this document with JS's printed revelations in 1833, and when it was included in the Doctrine and Covenants in 1835, no date was mentioned. Years later, the editors of JS's history mistakenly wrote that JS dictated this document in January 1832 in conjunction with his revision of the New Testament. But the location, which Whitmer identified as Wayne County, New York, indicates that the document was created before JS moved to Ohio in the second half of January 1831, and that it could have been written as early as January 1830.''%
				\footnote{%
					% \label{}%
					JSP footnote: ``JS could have been in Wayne County at almost any point during 1830, but there is some likelihood this document was produced during April, when JS was known to be in Wayne County. During that month, JS began baptizing believers and organizing branches of the church in New York; he also dictated two revelations that gave directions regarding baptism. Conversations with new and prospective converts may also have led to discussion about scriptural passages regarding infant baptism. In early December 1830, JS met recent convert Sidney Rigdon, who opposed infant baptism. JS and Rigdon worked together in December on JS's Bible revision, and during that period JS dictated a passage about baptism, but at that time JS and Rigdon were working in Seneca County. Similarly, any discussions JS might have had with James Covel would likely have taken place in Ontario or Seneca counties, not in Wayne County. If this document did result from contact with Covel, it could possibly have been produced in late January 1831 when JS and Rigdon traveled through Wayne County on their way to Ohio.''
					}
		
	% -----
	\subsubsection{Versions}
	% -----
		\label{DC74-V}
		
		See the \href{https://drive.google.com/file/d/1goAvNz8jS4R8slYfMG_db55Ty2z_vmgK/view?usp=sharing}{sections comparison} document.
	
		\begin{compactdesc}
			\item[JSP] \hyperref[EMS-1830-JS-1830]{1830}
			\item[BC] ---
			\item[DC 1835] Section 73, 202--203
			\item[TS] 5:13 (15 July 1844), 576
			\item[DC 1844] Section 74, 309--310
			\item[MS] 14:8 (15 April 1852), 116
		\end{compactdesc}

	% -----
	\subsubsection{Section 74:1} % *DC74-1 
	% -----
		\label{DC74-1}

			\footnote{%
				% \label{}%
				This first verse is \hyperref[1CORINTHIANS-7-14]{1 Corinthians 7:14}.
				}%
		FOR the unbelieving husband is sanctified by the wife, and the unbelieving wife is sanctified by the husband; else were your children unclean, but now are they holy.

		% \begin{framed}
		% \scriptsize
			% \begin{compactdesc}
				% \item[JSP] 
				% % \item[BC] 
				% % \item[]
			% \end{compactdesc}
		% \end{framed}

	% -----
	\subsubsection{Section 74:2} % *DC74-2 
	% -----
		\label{DC74-2}

			\footnote{%
				% \label{}%
				Here begins the explanation of the verse.
				}%
		Now, in the days of the apostles the law of circumcision was had among all the Jews who believed not the gospel of Jesus Christ.

		% \begin{framed}
		% \scriptsize
			% \begin{compactdesc}
				% \item[JSP] 
				% % \item[BC] 
				% % \item[]
			% \end{compactdesc}
		% \end{framed}

	% -----
	\subsubsection{Section 74:3} % *DC74-3 
	% -----
		\label{DC74-3}

		And it came to pass that there arose a great contention among the people concerning the law of circumcision, for the unbelieving husband was desirous that his children should be circumcised and become subject to the law of Moses, which law was fulfilled.

		% \begin{framed}
		% \scriptsize
			% \begin{compactdesc}
				% \item[JSP] 
				% % \item[BC] 
				% % \item[]
			% \end{compactdesc}
		% \end{framed}

	% -----
	\subsubsection{Section 74:4} % *DC74-4 
	% -----
		\label{DC74-4}

		And it came to pass that the children, being brought up in subjection to the law of Moses, gave heed to the traditions of their fathers and believed not the gospel of Christ, wherein they became unholy.

		% \begin{framed}
		% \scriptsize
			% \begin{compactdesc}
				% \item[JSP] 
				% % \item[BC] 
				% % \item[]
			% \end{compactdesc}
		% \end{framed}

	% -----
	\subsubsection{Section 74:5} % *DC74-5 
	% -----
		\label{DC74-5}

		Wherefore, for this cause the apostle wrote unto the church, giving unto them a commandment, not of the Lord, but of himself, that a believer should not be united to an unbeliever; except the law of Moses should be done away among them,

		% \begin{framed}
		% \scriptsize
			% \begin{compactdesc}
				% \item[JSP] 
				% % \item[BC] 
				% % \item[]
			% \end{compactdesc}
		% \end{framed}

	% -----
	\subsubsection{Section 74:6} % *DC74-6 
	% -----
		\label{DC74-6}

		That their children might remain without circumcision; and that the tradition might be done away, which saith that little children are unholy; for it was had among the Jews;

		% \begin{framed}
		% \scriptsize
			% \begin{compactdesc}
				% \item[JSP] 
				% % \item[BC] 
				% % \item[]
			% \end{compactdesc}
		% \end{framed}

	% -----
	\subsubsection{Section 74:7} % *DC74-7 
	% -----
		\label{DC74-7}

		But little children are holy, being \hyperref[SANCTIFICATION]{sanctified} through the atonement of Jesus Christ; and this is what the scriptures mean.

		% \begin{framed}
		% \scriptsize
			% \begin{compactdesc}
				% \item[JSP] 
				% % \item[BC] 
				% % \item[]
			% \end{compactdesc}
		% \end{framed}